\chapter{Experiments}

\section{Introduction}
This chapter details the way in which the experimentation has been conducted. While chapter~4 described both methodology and an evaluation framework, this chapter concerns itself with the actual process of conducting the experiment (software setup, models and services startup, execution of the notebooks, the sequence of scenario questions to be answered, parameterization of all aspects of the experiments, logging results and validating all outputs generated by each experiment).

While the organization of the experiments is based upon the local execution phase, there will also be an equivalent organization for the HPC phase, however, at this time the implementation is focused upon local comparisons and as such the majority of information contained within these result files relates to that aspect. It is necessary to note that it is the local executions which serve as the empirical baseline, whereas the HPC phase should execute the same notebooks and scenarios under identical conditions (i.e., without changing the retrieval logic).

\section{Experimental Setup}

\subsection{Software Environment}
The experiments were implemented as a series of Python Jupyter Notebooks. Each experiment was set up to use a Python Virtual Environment. All dependencies required by the repository were installed. Once the dependencies were installed all the experimentation Notebooks were loaded and run directly from the Project Notebooks Directory. The Local Workflow utilized Ollama for Model Serving and Milvus for both Vector-Database and Hybrid-Retrieval Variants of RAG.

The executed implementation notebooks are:

\begin{itemize}
    \item \texttt{rag\_\allowbreak basic.ipynb};
    \item \texttt{rag\_\allowbreak vectordb.ipynb};
    \item \texttt{rag\_\allowbreak vectordb\_\allowbreak tuned.ipynb};
    \item \texttt{rag\_\allowbreak reranker.ipynb};
    \item \texttt{rag\_\allowbreak hybrid.ipynb}; and
    \item \texttt{rag\_\allowbreak hybrid\_\&\_\allowbreak rerank.ipynb}.
\end{itemize}

All notebooks followed a similar general process of execution. That process included importing libraries, configuring paths and services, loading the document corpus, splitting/indexing documents, building the RAG pipeline, executing scenario questions, and saving results. This standardized structure allows for the elimination of unwanted variation between implementations.

\subsection{Models and Services}
In order to run locally two Ollama models need to exist. The generation model will be \texttt{llama3.1} and the embedding model will be \texttt{mxbai-embed-large}. For dense retrieval, vector-database retrieval, and hybrid retrieval, all of these methods rely upon the embeddings. Reranking notebooks also make use of a MiniLM cross-encoder from the sentence-transformers ecosystem.

Prior to running the vector-databases (reranking), hybrid, tuned-vector-databases, hybrid-rerank notebooks, the milvus service is launched locally. A milvus local service is managed by the project script \texttt{standalone\_\allowbreak embed.sh}. An in-memory baseline can be used without having milvus installed and thus can be used for a reference implementation when you want something simple.

At the HPC stage later in this paper, we wish to maintain as many of the model choices and retrieval configurations as possible. In cases where changes are necessary due to differences in serving methods, container images, file system paths, or scheduler configurations those items should have been noted prior to comparing local versus HPC results.

\subsection{Input Corpus}
All experiments make use of the same OrchestrAI synthetic enterprise corpus. The corpus includes 54 DOCX documents divided by departments. Each notebook will load each document recursively and convert them into a haystack document format; it also keeps track of source information about each document (file name and department).

Since all notebooks have access to the same set of documents, the primary cause for differences in answers is the way in which they retrieve and rank those documents. The same corpus needs to be duplicated on the HPC system used for the second phase and retain the directory structure and file names.

\subsection{Configuration and Parameters}
The key configurable options include the model types, query types (e.g. "what," "where"), total number of retrieved documents, chunking configurations and logging schema. In most cases we have passed in a total of ten retrieved document-chunks to the language models for evaluation. For the reranked versions we first retrieve twenty candidates using our ranking models, and then filter down to the top ten based on the output of our cross encoder rerankers.

All of the chunking configurations differ from each other only when they were used as part of an experimental design. The baseline used 120 word chunks that had 20 words of overlap. All of the reranking variants including the dense milvus version, the hybrid version and the hybrid-reranking variant all used 80 word chunks with 20 words of overlap. The only chunk size used by the vector database tuning notebook was 200 words with 40 words of overlap. The rest of the vector-database configuration remained unchanged.

\begin{itemize}
    \item \textbf{Top 10 retrieval:} basic, vector database, tuned vector database, and hybrid.
    \item \textbf{Top 20 retrieval followed by top 10 reranking:} reranking and hybrid with reranking.
    \item \textbf{Chunking test:} tuned vector database changes only chunk length and overlap.
\end{itemize}

\section{Deployment Stages}

\subsection{Local Execution Stage}
The local execution stage represents the final experimental stage of the project. The primary goal of this stage is to validate that each notebook successfully runs on the author’s computer and to provide comparative timing, retrieval and answer quality outputs.

The notebooks were run in the same order as follows:

\begin{enumerate}
    \item baseline in-memory \ac{RAG};
    \item Milvus vector-database \ac{RAG};
    \item tuned Milvus vector-database \ac{RAG};
    \item dense retrieval with reranking;
    \item hybrid retrieval; and
    \item hybrid retrieval with reranking.
\end{enumerate}

The above order represented an increasing complexity pipeline, starting with the most basic pipeline and adding persistent vector storage, chunking variations, reranking, hybrid retrieval and finally hybrid retrieval with reranking. All outputs generated during each run are recorded into a single results.CSV file for comparison purposes.

\subsection{HPC Execution Stage}
% TODO: Complete this subsection after the HPC experiments are executed.
% TODO: Add job scripts, environment setup, node details, scheduler settings,
% storage paths, and container/service startup details.

\section{Experiment Scenarios}

\subsection{Scenario-Based Query Execution}
Five predetermined question-answering scenarios serve as the framework for the experiments. Different retrieval and reasoning behaviours are tested in each case. A straight factual lookup is the first scenario. The latter situations call for cross-document or procedural reasoning.

The five scenario questions are:

\begin{enumerate}
    \item \textbf{\textit{What is the first-response SLA for a Sev-1 support incident?}} This is the clean baseline scenario. It checks whether the system can retrieve and answer a single operational fact accurately.
        \begin{itemize}
            \item Retrieval type: single-document
            \item Main test: exact fact retrieval
        \end{itemize}
    \item \textbf{\textit{What steps must be completed before a new employee is considered ready for day one at OrchestrAI?}} This scenario is slightly more complex because the answer is not a single fact but a short operational process.
        \begin{itemize}
            \item Retrieval type: checklist or procedure retrieval
            \item Main test: whether the system can recover and summarize a structured process
        \end{itemize}
     \item \textbf{\textit{A new engineer starts on Monday, but their laptop has not arrived and their required access is still pending. Which teams are involved and what should happen next?}} This scenario requires combining information from more than one team to answer a realistic workplace problem.
        \begin{itemize}
            \item Retrieval type: multi-document, cross-team retrieval
            \item Main test: whether the system can synthesize HR, IT, and access-related responsibilities
        \end{itemize}
     \item \textbf{\textit{Can OrchestrAI make an urgent hardware purchase without following the normal PO workflow if a warehouse outage is at risk?}} This is the most demanding scenario in the notebook because it requires cross-functional reasoning across several business areas.
        \begin{itemize}
            \item Retrieval type: policy retrieval with exception logic
            \item Main test: whether the system can identify when a standard workflow may be bypassed
        \end{itemize}
     \item \textbf{\textit{A customer renewal is at risk because support cases are recurring, a requested feature is still unavailable, and replacement hardware is delayed. What should OrchestrAI do next, and which teams must be involved?}} This scenario is slightly more complex because the answer is not a single fact but a short operational process.
        \begin{itemize}
            \item Retrieval type: multi-hop, multi-document retrieval
            \item Main test: whether the system can combine support, product, logistics, and commercial risk signals into one grounded answer
        \end{itemize}
\end{enumerate}

The same questions are used for all implementations. Each implementation uses its own scenario identifier prefix so that rows remain easy to distinguish in the results file. For example, the baseline uses \texttt{b1}--\texttt{b5}, the vector-database run uses \texttt{v1}--\texttt{v5}, and the hybrid-reranking run uses \texttt{hr1}--\texttt{hr5}.

\subsection{Baseline RAG Experiment}
In this baseline experiment we execute the five scenarios with our in memory \ac{RAG} pipeline. We load the DOCX chunks in memory and index them into an in memory document store. For each of the five questions we retrieve the top 10 dense match candidates. This experiment will serve as a reference point for all other experiments presented in this paper. In addition it serves as a baseline for the other experiments.

\subsection{Vector Database RAG Experiment}
This experiment executes the same five scenarios, but now on a persistent vector database (Milvus). Again, we search for the top 10 dense match candidates for each of the five questions. This experiment examines the impact of persistently storing vectors in a database and if the quality of the retrieved matches is still similar to those produced by the baseline experiment.

\subsection{Tuned - Vector Database RAG Experiment}
Additionally, this experiment also includes the execution of the same five scenarios with a larger chunk size and overlap. As previously mentioned, chunk size has a strong influence on the behavior of the retriever and the answers generated. To isolate the effects of chunk size and maintain consistent conditions for the retriever, model, prompt, and top-\( k \) value we use the same setting for the retriever, model, prompt, and \( k \) value. Therefore we call this experiment "tuned".

\subsection{Reranking Experiment}
For this reranking experiment, we again execute the five scenarios based on dense retrieval (Milvus) but now followed by cross encoder reranking. First, the retriever delivers 20 candidates and then the reranker selects the best 10 candidates for prompting. This experiment assesses the trade-off between additional latency during retrieval stage versus improved ranking order.

\subsection{Hybrid Retrieval Experiment}
The hybrid retrieval experiment combines dense semantic retrieval with a sparse BM25-style signal generated inside Milvus. It retrieves 10 documents for each question. This experiment evaluates whether combining semantic and lexical signals improves robustness for questions that contain exact operational terms such as severity labels, team names, procurement terms, and workflow names.

\subsection{Hybrid Retrieval with Reranking Experiment}
In the hybrid reranking experiment, we apply cross encoder reranking after hybrid candidate retrieval. The hybrid retriever provides 20 candidates and the reranker chooses the best 10. This is the most complex pipeline currently implemented. This experiment tests whether hybrid recall combined with reranked precision generates better answers than previous implementations at a higher computational cost.

\section{Benchmarking and Data Collection}
All results are stored in a shared CSV file via our project's results logger. Prior to running each notebook, we clear out any rows for that specific implementation so repeated runs overwrite old data rather than produce multiple entries. Each row includes implementation name, scenario ID, question, generated answer, source documents, departmental information from retrieved sources, context information from retrieved sources, number of retrieved documents, top-\( k \) settings, indexing time, retrieval time and generation time.

It is crucial to note that there is an important difference between indexing time, which represents one-time costs associated with initial document preparation and storage; whereas retrieval time and generation time represent ongoing costs during answer production.

Once all scenario runs have completed, we take the results file to create the manual evaluation worksheet. Using evidence from the DOCX corpus we fill in manual scores and supporting documents. Then our deterministic evaluation script calculates several relevant metrics including retrieval and answer quality and creates two output files per-scenario and summary CSV files. Our visualization notebook then utilizes these preprocessed outputs to create various types of plots including answer-quality, retrieval-quality, scenario-based plots, timing-based plots, and radar-charts for Chapter 6.

% TODO: Add environment field before combining local and HPC results.

\section{Data Validation}
The recorded output is validated as complete prior to analysis. The validation entails verification that there are five rows in each of the implementations, that the generated answers were not null, that there were sources that were retrieved, and that the time fields were available. Additionally, all retrieved contexts are manually examined when developing the truth table (ground-truth sheet), since the quality score of the answer cannot be assumed; it must have some form of evidence based upon a subset of the data within the corpus.

% TODO: Complete HPC validation after HPC outputs are generated.

\section{Summary}
This chapter describes the experimental setting and execution for the developed \ac{RAG} pipelines. It provides information on the software environment, model services, input corpus, notebooks execution order, deployment phases, question types used for scenarios, specific experiment designs for each of the pipelines, log results for each of the pipelines, and validation procedure. The local phase is where the comparison took place for the different retrieval approaches with respect to their performance by comparing them on a local computer. 

% TODO: Complete HPC say about the same tests on HPC
